\ncdlesson{4}{Kunsten at være uenig på en høflig måde}{礼貌表达不同意见的艺术}
{DU3 Modul 4（B1+ $\rightarrow$ B2-）｜对标 NCE 3 \& Modultest 4 \textit{Skrivning Opgave 2 (Debatindlæg) \& Mundtlig Opgave 2 (Paropgave)}}
{\textbf{让步与反驳高阶句式}（\textit{Ganske vist..., men ikke desto mindre... / Det kan godt være, at..., men... / Selvom..., mener jeg dog...}）· 让步连接词与副词的\textbf{主从句语序与 V2 倒装} · 丹麦式“先认可、再温和反驳”的共识协商修辞}
{在丹麦职场讨论与 DU3 Modul 4 口试、写作中，直接甩出一句 \textit{”Det er jeg helt uenig i!”} 或 \textit{”Du tager fejl!”} 往往显得生硬且具攻击性。丹麦社会高度重视\textbf{共识文化（\textit{konsensuskultur}）}与平实客气的讨论氛围：高分表达讲究“先退一步承认对方合理之处，再稳健有力地提出保留或反驳”。吃透本课的 \textbf{\textit{Ganske vist..., men ikke desto mindre...}} 与 \textbf{\textit{Det kan godt være, at..., men...}}，你的不同意见就能既有礼貌又有锋芒！}

% ---------------------------------------------------------------------
\begin{foerduleaeser}
带着下面两个问题阅读 Tekst A 和 Tekst B：
\begin{enumerate}
  \item \textit{Hvorfor oplevede Marco det først som forvirrende, når hans danske kolleger var uenige på møder i Tekst A, og hvordan formulerer de i stedet saglig kritik?}（在 Tekst A 中，为什么马尔科最初觉得丹麦同事在会上表达不同意见的方式令人困惑？他们实际是如何提出客观批评的？）
  \item \textit{Hvad er Henrik og Sofie henholdsvis enige og uenige om i debatten om hjemmearbejdsdage i Tekst B?}（在 Tekst B 关于居家办公日的讨论中，亨里克和索菲各自赞同什么、又在哪些方面存在分歧？）
\end{enumerate}
\end{foerduleaeser}

% ---------------------------------------------------------------------
\ncdsubsec{二、Tekster（口语微调 / 报刊与思辨双语篇）}

\subsection*{Tekst A：\textit{Når et ”nej” pakkes ind i høflighed}（职场文化随笔 · 约 220 词）}

\begin{tipbox}{语篇说明}
本篇以意大利籍工程师马尔科（Marco）初入丹麦职场的经历切入，剖析丹麦人如何在保留体面与合作氛围的同时表达反对意见。阅读时请重点标出\textbf{\color{NCDNavy}深蓝色加粗的让步引导语}（先认可对方部分道理）与\textbf{\color{NCDRed}深红色加粗的反驳转折接口}（注意 \textit{Ganske vist} 与 \textit{ikke desto mindre} 触发的 V2 倒装！）。
\end{tipbox}

Da Marco begyndte som ingeniør i en virksomhed i Odense, troede han i lang tid, at alle på afdelingen var helt enige om næsten alt. På projektmøderne hørte han nemlig sjældent nogen sige direkte, at et forslag var dårligt. Først efter nogle måneder opdagede han, at hans danske kolleger ofte udtrykte skarp uenighed gennem faste høflige mønstre.

Når en kollega for eksempel foreslog en billig løsning på et teknisk problem, svarede Emma roligt: ”\da{Det kan godt være, at} den model er billigere her og nu, \dabold{men} på lang sigt kræver den langt mere vedligeholdelse.” Dermed anerkendte hun først kollegaens pointe, før hun venligt pegede på svagheden.

Senere lærte Marco, at man i skriftlige vurderinger og formelle debatter ofte bruger endnu stærkere strukturer til at skabe balance. En teamleder forklarede ham princippet med et klassisk eksempel: ”\da{Ganske vist} sparer vi penge i år, \dabold{men ikke desto mindre} risikerer vi at miste vigtige kunder næste år.” Her viser \textit{ganske vist}, at man tager modpartens argument alvorligt, mens \textit{ikke desto mindre} fastholder den afgørende indvending.

I dansk arbejdskultur handler det nemlig om at adskille person og sag. \da{Selvom} man kan være dybt uenig i en konklusion, \dabold{bør} tonen \dabold{dog} forblive respektfuld og saglig. Når man mestrer denne balance mellem indrømmelse og modargument, bliver man taget langt mere alvorligt end den, der blot råber højest.

\begin{transbox}
\small
\textbf{\color{NCDTeal}[Tekst A 中文精译]}\\
当马尔科开始在欧登塞的一家企业担任工程师时，有很长一段时间他都以为部门里所有人在几乎所有事情上都完全一致。因为在项目会议上，他很少听到有人直接说某个方案很糟糕。直到几个月后他才发现，丹麦同事其实常常通过固定的礼貌句式来表达鲜明的反对意见。

例如，当一位同事针对某个技术难题提出了一个廉价方案时，埃玛平静地回应道：“那个模型眼下确实可能更便宜，但从长远来看，它需要多得多的维护成本。”她借此先认可了同事的着眼点，然后才温和地指出其中的薄弱之处。

后来马尔科了解到，在书面评估和正式辩论中，人们还常常使用更有力度的结构来营造平衡感。一位团队负责人用一个经典例句向他解释了这一原则：“我们今年诚然省下了资金，然而尽管如此，明年我们却面临着失去重要客户的风险。”在这里，\textit{ganske vist} 表明自己认真对待了对方的论据，而 \textit{ikke desto mindre} 则坚定地守住了关键的反驳立场。

因为在丹麦职场文化中，关键在于把“人”和“事”分开。尽管大家在某项结论上可能存在深刻分歧，语气却仍然应当保持尊重且就事论事。一旦掌握了这种在让步与反论之间游刃有余的平衡艺术，你的意见就会比一味嗓门最大的人受到重视得多。
\end{transbox}

\subsection*{Tekst B：\textit{Debat på intranettet: Skal alle tilbage på kontoret?}（企业内网正反讨论 · 约 125 词）}

\begin{emailbox}{职场论坛讨论｜Henrik (Projektleder) \& Sofie (Analytiker)}
\textbf{Henrik:} Kære kolleger. Jeg synes, vi bør reducere antallet af hjemmearbejdsdage fra to til én om ugen. \da{Det kan godt være, at} mange arbejder koncentreret hjemme ved skrivebordet, \dabold{men} vi mister alt for meget af den spontane sparring ved kaffemaskinen. \da{Ganske vist} fungerer onlinemøder fint til faste opgaver, \dabold{men ikke desto mindre} svækkes vores fællesskab i hverdagen.

\textbf{Sofie:} Tak for dit oplæg, Henrik. Jeg forstår godt din bekymring for holdånden, og jeg er enig med dig i, at fysiske møder er værdifulde. \da{Selvom} uformelle snakke på gangen er vigtige, \dabold{mener} jeg \dabold{dog}, at to hjemmedage stadig bør bevares. \da{Ganske vist} er vi færre på kontoret om fredagen, \dabold{men til gengæld} sparer mange pendlere både tid og stress, hvilket gør os mere effektive resten af ugen.
\end{emailbox}

\textbf{Tekst B 阅读核对（Rigtigt / Forkert / Kort svar，答案见第六部分）：}
\begin{enumerate}
  \item \textit{Henrik mener slet ikke, at onlinemøder fungerer til nogen arbejdsopgaver.} \underline{\hspace{2.2cm}}
  \item \textit{Hvad er Henrik bekymret for, at virksomheden mister med to hjemmearbejdsdage?} \underline{\hspace{2.2cm}}
  \item \textit{Hvad peger Sofie på som den konkrete fordel ved, at medarbejdere arbejder hjemme to dage om ugen?} \underline{\hspace{2.2cm}}
\end{enumerate}

\begin{tipbox}{本课发音与语调要点}
\begin{itemize}
  \item \textbf{高阶让步连词与副词读音}：\textit{ganske vist} \ipa{ˈɡansɡə ˈvesd}（\textit{ganske} 词尾轻读 \ipa{ɡə}）；\textit{ikke desto mindre} \ipa{ˈeɡə ˈdɛsdo ˈmenˀdʁə}（重音落在 \textit{ikke}、\textit{des-} 与 \textit{min-} 上，带喉塞音）；\textit{alligevel} \ipa{ɑˈliːəˌvɛlˀ}（中间的 \textit{g} 完全不发音！）；\textit{selvom} \ipa{ˈsɛlˀˌʌm}（注意词尾 \textit{v} 在口语中常轻读或接近脱落）。
  \item \textbf{立场与争论高频词的软 \textit{d} 与静音字母}：\textit{uenig} \ipa{uˈeːˀni}（前缀 \textit{u-} 重读或次重读，词尾 \textit{g} 不发音！）；\textit{vedligeholdelse} \ipa{veðˈliːəˌhʌlˀəlsə}（第一个 \textit{d} 为软 \textit{d}，\textit{g} 与第二个 \textit{d} 均不发音）；\textit{at anerkende} \ipa{ˈanɐˌkʰɛnˀə}（\textit{-nd-} 中的 \textit{d} 不发音！）；\textit{saglig} \ipa{ˈsɑʊ̯li}（\textit{g} 滑音化，词尾 \textit{g} 不发音）。
  \item \textbf{礼貌语调提示}：说 \textit{Det kan godt være, at...} 时，重音轻轻放在 \textit{godt} 上，尾音微微下沉留顿，再平稳接出 \textit{men...}，听起来诚恳而不带挑衅感。
\end{itemize}
\end{tipbox}

\begin{transbox}
\small
\textbf{\color{NCDTeal}[Tekst B 中文对照]}\\
\textbf{亨里克（项目经理）}：各位同事好。我认为我们应当把每周居家办公的天数从两天减到一天。诚然，许多人在家里的书桌前能非常专注地工作，但我们在茶水间丧失了太多自发的灵感碰撞。虽然线上会议对于处理固定任务确实运作良好，然而尽管如此，我们日常的团队凝聚力却正在被削弱。

\textbf{索菲（分析师）}：谢谢你的提议，亨里克。我完全理解你对团队精神的担忧，我也同意你所说的面对面会议很有价值。尽管走廊里的非正式交谈很重要，但我依然认为应当保留两天居家办公。诚然我们周五在办公室的人少了一些，但作为回报，许多通勤员工既省下了时间又减轻了压力，这使得大家在一周其余时间里效率更高。
\end{transbox}

% ---------------------------------------------------------------------
\ncdsubsec{三、Gloser（核心生词表）}

\begin{tipbox}{词表重点}
表达“同意/不同意”时，介词搭配非常固定：和某人意见一致/不合用 \textit{med}（\textit{uenig med Henrik}），在具体观点上同意/不同意用 \textit{i}（\textit{uenig i forslaget / uenig i, at...}）。此外，名词化词 \textit{en vedligeholdelse}（维护）和 \textit{en anerkendelse}（认可）均为通性 \textit{en} 名词，动词 \textit{at modsige}（反驳）为不规则动词（随 \textit{sige} 变化为 \textit{modsagde, har modsagt}）。
\end{tipbox}

\begin{glosertable}
副词短语 & ganske vist & \textit{(不变，让步副词短语)} & 诚然，固然，确实如此 \\
副词短语 & ikke desto mindre & \textit{(不变，转折副词短语)} & 尽管如此，然而依然 \\
连词 & selvom & \textit{(不变，让步从属连词)} & 尽管，虽然（后接从句语序） \\
副词 & alligevel & \textit{(不变，让步副词)} & 尽管如此，还是，仍然 \\
副词 & dog & \textit{(不变，句中转折副词)} & 不过，然而（多置于动词后） \\
\midrule
名词 & en uenighed & \textit{(uenigheden, uenigheder, -erne)} & 分歧，意见不合 \\
名词 & en konsensus & \textit{(konsensussen, ingen flertal)} & 共识，一致意见 \\
名词 & en pointe & \textit{(pointen, pointer, pointerne)} & 要点，核心论点 \\
名词 & et modargument & \textit{(modargumentet, -argumenter, -erne)} & 反论据，反驳理由 \\
名词 & en sparring & \textit{(sparringen, ingen flertal)} & 业务切磋，思想碰撞 \\
名词 & en vedligeholdelse & \textit{(vedligeholdelsen, ingen flertal)} & 维护，保养 \\
名词 & en høflighed & \textit{(høfligheden, ingen flertal)} & 礼貌，客气 \\
\midrule
动词 [har] & at anerkende & \textit{(anerkender, anerkendte, har anerkendt)} & 认可，承认（对方合理之处） \\
动词 [har] & at indrømme & \textit{(indrømmer, indrømmede, har indrømmet)} & 承认，作出让步 \\
动词 [har] & at modsige & \textit{(modsiger, modsagde, har modsagt)} & 反驳，反口，相矛盾 \\
动词 [har] & at fastholde & \textit{(fastholder, fastholdt, har fastholdt)} & 坚持（观点、原则） \\
\midrule
形容词 & uenig & \textit{(uenigt, uenige)} & 不同意的（与某人 \textit{med} / 对某事 \textit{i}） \\
形容词 & saglig & \textit{(sagligt, saglige)} & 客观务实的，就事论事的 \\
\end{glosertable}

% ---------------------------------------------------------------------
\ncdsubsec{四、Grammatik \& Stilistik（核心语法、语气与修辞精讲）}

\begin{grammatikbox}{语法点 1：让步与反驳四大高阶句式与语序全景对照}
在 DU3 Modul 4 写作（尤其是 \textit{Debatindlæg}）和口语双人辩论中，“先让步再反驳”是最具说服力的修辞骨架。请严格区分以下四种句式的\textbf{正式度}与\textbf{V2 倒装位置}：

\vspace{0.3em}
\renewcommand{\arraystretch}{1.24}
{\small
\begin{tabularx}{\linewidth}{@{} >{\bfseries\raggedright\arraybackslash}p{3.1cm} >{\raggedright\arraybackslash}p{4.3cm} >{\raggedright\arraybackslash}X @{}}
\toprule
\textbf{让步反驳句式} & \textbf{语域与语序铁律} & \textbf{标准政论与协商例句} \\
\midrule
\da{Ganske vist...,}\newline \da{men ikke desto mindre...} &
\textbf{书面正式高阶}（报刊/辩论首选）\newline
\textit{Ganske vist} 放句首必 \textbf{V2 倒装}；\textit{men} 后紧接 \textit{ikke desto mindre} 也必 \textbf{V2 倒装}！ &
\textit{\da{Ganske vist} \dabold{sparer vi} penge nu, \da{men ikke desto mindre} \dabold{risikerer vi} store fejl senere.}\newline
（也可分两句：\textit{... \da{Ikke desto mindre} \dabold{mangler} vi tal.}） \\
\midrule
\da{Det kan godt være, at..., men...} &
\textbf{口语与半正式通用}（委婉留余地）\newline
\textit{at} 后守\textbf{从句语序}（\textit{ikke} 前移）；\textit{men} 后接正常主句 &
\textit{\da{Det kan godt være, at} planen \dabold{ikke er} perfekt, \da{men} den løser problemet hurtigt.} \\
\midrule
\da{Selvom...,}\newline \da{... dog / alligevel...} &
\textbf{标准主从让步句}\newline
\textit{Selvom} 内部不倒装；前置时后半主句必须 \textbf{V2 倒装}，且\textbf{绝不加 \textit{men}}！ &
\textit{\da{Selvom} møder tager tid, \dabold{mener} jeg \da{dog}, at dialogen er vigtig.}\newline
（\textit{dog} 放在主句动词或主语后） \\
\midrule
\da{Jeg forstår godt..., men...}\newline \da{Jeg er kun delvist enig i, at...} &
\textbf{高情商开篇表态句}\newline
专用于回应对方来信或辩论贴的第一句，立刻奠定客观礼貌基调 &
\textit{\da{Jeg er kun delvist enig i}, at alle bør møde på kontoret hver dag.} \\
\bottomrule
\end{tabularx}
}
\vspace{0.3em}

\textbf{语序灵活微调提示}：\ding{172} \textit{ganske vist} 和 \textit{ikke desto mindre} 也可以从句首移到主句动词（或主语）后面，以让语气更紧凑自然：\textit{Forslaget er \da{ganske vist} dyrt, men det har \da{ikke desto mindre} stor værdi.} \ding{173} 想在反驳句中突出“作为补偿/相反另一方面”，还可用 \da{men til gengæld...} 或 \da{men på den anden side...}（放句首或 \textit{men} 后均须 V2 倒装）。
\end{grammatikbox}

\begin{grammatikbox}{语法点 2：中英对比排雷——让步反驳最易踩中的四大扣分陷阱}
\begin{enumerate}
  \item \textbf{严禁把“虽然……但是……”直译成 \textit{Selvom..., men...}！}：这是中文母语者最顽固的母语迁移错误。在丹麦语中，\textit{Selvom}（从属连词）和 \textit{men}（并列连词）\textbf{不能同时出现在同一主从句里}！\newline
    \ding{55}\ \textit{Selvom det er dyrt, men vi køber det alligevel.}\newline
    \ding{51}\ \textit{Selvom det er dyrt, \dabold{køber} vi det \da{alligevel} / \da{dog}.}（后半主句直接动词倒装！）
  \item \textbf{\textit{Ganske vist} 与 \textit{Ikke desto mindre} 放句首绝不打英语逗号！}：英语习惯写 \textit{Admittedly, ... / Nevertheless, ...}；丹麦语绝不加逗号，且动词必须紧随其后占据第 2 位：\ding{51}\ \textit{\da{Ganske vist} \dabold{er} mødet langt, men \da{ikke desto mindre} \dabold{bør} vi deltage.}（绝对不能写成 \ding{55}\ \textit{Ganske vist, mødet er langt...}）。
  \item \textbf{记住 \textit{Det kan godt være, at...} 从句内的否定词位置}：在 \textit{at} 引导的从句里，\textit{ikke / aldrig} 必须跑到变位动词前面：\ding{51}\ \textit{Det kan godt være, at han \dabold{ikke har} tid i dag, men...}（不能写成 \ding{55}\ \textit{...at han har ikke tid}）。
  \item \textbf{分清 \textit{enig / uenig med}（对人）与 \textit{enig / uenig i}（对观点/从句）}：\newline
    \ding{51}\ \textit{Jeg er uenig \dabold{med} dig \dabold{i}, at forbuddet hjælper.}（我和你在“禁令有效”这件事上意见不合——先 \textit{med} 人，再 \textit{i} 事！）。
\end{enumerate}
\end{grammatikbox}

% ---------------------------------------------------------------------
\ncdsubsec{五、Mønsterøvelser（句型与修辞操练）}

\textbf{练习 A：高阶让步反驳精准填空（每空填入唯一准确的 1 个丹麦语词）}
\begin{enumerate}
  \item \underline{\hspace{2.2cm}} vist er hjemmearbejde meget fleksibelt, men vi mister noget af den spontane sparring.
  \item Forslaget lyder spændende, men ikke \underline{\hspace{2.2cm}} mindre må vi først vurdere de samlede omkostninger.
  \item Jeg er desværre uenig \underline{\hspace{2.2cm}} min chef i, at alle møder nødvendigvis skal være fysiske.
  \item Vi er kun delvist enige \underline{\hspace{2.2cm}}, at den nye arbejdsform styrker trivslen for alle medarbejdere.
  \item Det kan godt \underline{\hspace{2.2cm}}, at løsningen tager lidt tid at indføre, men den holder til gengæld længe.
  \item \underline{\hspace{2.2cm}} mange foretrækker at arbejde hjemme om fredagen, mener ledelsen dog, at kontoret er vigtigt.
  \item Argumentet har en god pointe. Ikke desto mindre \underline{\hspace{2.2cm}} planen stadig flere alvorlige svagheder. (rumme，现在时)
\end{enumerate}

\textbf{练习 B：让步反驳句型升级与语病修正（按括号要求改写，或修正中/英式常见错误）}
\begin{enumerate}
  \item 用 \textit{Ganske vist..., men ikke desto mindre...} 合并两句：\textit{Kaffemaskinen er dyr. Den skaber et bedre arbejdsmiljø.} $\rightarrow$
  \item 消除“虽然……但是”叠加病句：\textit{Selvom projektet kræver mange timer, men jeg synes, vi bør gennemføre det.} $\rightarrow$
  \item 用 \textit{Det kan godt være, at..., men...} 礼貌反驳下方生硬句子：\textit{Din idé er dårlig, for den koster alt for meget.}（先承认“这主意有创意”）$\rightarrow$
  \item 修正句首逗号与倒装错误：\textit{Ganske vist, planen har sine fordele. Ikke desto mindre, vi bør vente med beslutningen.} $\rightarrow$
  \item 修正从句语序与介词搭配：\textit{Jeg er uenig på dig om, at man kan ikke arbejde effektivt hjemmefra.} $\rightarrow$
  \item 用 \textit{Jeg er kun delvist enig i, at...} 温和改写：\textit{Jeg er helt uenig med, at onlinemøder altid er spild af tid.} $\rightarrow$
\end{enumerate}

\textbf{练习 C：中译丹实战（活用礼貌让步反驳句式与精准语序）}
\begin{enumerate}
  \item 我部分同意你的观点，但我们依然需要更多数据。
  \item 这项计划诚然很昂贵，然而尽管如此它却极具长远价值。
  \item 也许他说得没错，不过我认为我们还是应该等待。
  \item 尽管大家意见不合，团队还是保持着礼貌而务实的语气。
  \item 在数字化方案这个问题上，我与同事的看法并不一致。
  \item 尽管线上会议省时，它们却无法完全替代面对面的交流。
\end{enumerate}

% ---------------------------------------------------------------------
\ncdsubsec{六、DU3 Modul 4 实战微任务与参考答案}

\begin{emailbox}{DU3 Modul 4 Skrivning Opgave 2 实战微任务（Debatindlæg · minimum 100 ord）}
\textbf{Situation}: I lokalavisen har en læser skrevet et skarpt debatindlæg om moderne kontorer:\\
\textit{”Åbne storrumskontorer bør afskaffes med det samme! Man kan slet ikke koncentrere sig på grund af støj fra kolleger og telefoner. Alle bør i stedet have deres eget lukkede kontor.”}\\[3pt]
\textbf{Opgave}: Skriv et høfligt og sagligt \textbf{debatindlæg} på \textbf{mindst 100 ord} (gerne 130--155 ord), hvor du forholder dig til indlægget. Anerkend først læserens pointe, forklar derefter, hvorfor et totalt forbud mod storrumskontorer ikke nødvendigvis er den bedste løsning, og brug mindst \textbf{to} indrømmelsesstrukturer fra denne lektion (\textit{Ganske vist..., men ikke desto mindre... / Det kan godt være, at..., men... / Selvom...}).
\end{emailbox}

\begin{answerbox}
\textbf{Før du læser 参考回答}
\begin{enumerate}
  \item \textit{Marco blev først forvirret, fordi han sjældent hørte nogen sige direkte på møderne, at et forslag var dårligt. Hans danske kolleger formulerer i stedet saglig kritik ved først at anerkende modpartens pointe med udtryk som ”Det kan godt være, at..., men...” eller ”Ganske vist..., men ikke desto mindre...”.}
  \item \textit{Henrik og Sofie er enige om, at fysiske møder er værdifulde. Men Henrik vil skære ned til én hjemmearbejdsdag om ugen, fordi han frygter, at man mister spontan sparring og fællesskab, mens Sofie mener, at to hjemmedage bør bevares, da pendlere sparer tid og stress og dermed arbejder mere effektivt.}
\end{enumerate}

\textbf{Tekst B 阅读核对答案}
\begin{enumerate}
  \item \da{Forkert}（亨里克明确承认线上会议用于固定任务效果很好：\textit{Ganske vist fungerer onlinemøder fint til faste opgaver}）
  \item \da{Den spontane sparring ved kaffemaskinen og fællesskabet i hverdagen.}（他担心失去茶水间自发的业务交流以及日常团队凝聚力）
  \item \da{Pendlere sparer både tid og stress, hvilket gør medarbejderne mere effektive resten af ugen.}（通勤员工节省时间和减轻压力，从而提高一周其余时间的效率）
\end{enumerate}

\textbf{练习 A 答案}
\begin{enumerate}
  \item \da{Ganske}（固定搭配：\textit{Ganske vist...}“诚然……”）
  \item \da{desto}（固定转折短语：\textit{ikke desto mindre}“然而尽管如此”）
  \item \da{med}（表示“与某人意见不合”必须用介词 \textit{med}：\textit{uenig med min chef}）
  \item \da{i}（在具体观点或从句前表示赞同/不赞同，固定用介词 \textit{i}）
  \item \da{være}（委婉让步开篇句式：\textit{Det kan godt være, at...}）
  \item \da{Selvom}（引导让步从句，注意后半主句不再加 \textit{men}，而用 \textit{dog}）
  \item \da{rummer}（\textit{Ikke desto mindre} 放句首触发 V2 倒装，紧接现在时动词 \textit{rummer}）
\end{enumerate}

\textbf{练习 B 答案（参考）}
\begin{enumerate}
  \item \textit{\da{Ganske vist} er kaffemaskinen dyr, \dabold{men ikke desto mindre} skaber den et bedre arbejdsmiljø.}（注意两处 V2 倒装：\textit{er kaffemaskinen} 与 \textit{skaber den}）
  \item \textit{\da{Selvom} projektet kræver mange timer, synes jeg \dabold{dog}, at vi bør gennemføre det.}（删掉多余的 \textit{men}，用倒装及句中 \textit{dog} 或 \textit{alligevel} 呼应）
  \item \textit{\da{Det kan godt være, at} din idé er kreativ, \dabold{men} den koster desværre alt for meget.}（先肯定对方创意，再委婉提出费用问题）
  \item \textit{\da{Ganske vist har} planen sine fordele. \dabold{Ikke desto mindre bør} vi vente med beslutningen.}（去掉英语式逗号，并保证句首副词后立刻接变位动词）
  \item \textit{Jeg er uenig \dabold{med} dig \dabold{i}, at man \dabold{ikke kan} arbejde effektivt hjemmefra.}（修正介词 \textit{med / i}，并将从句内的 \textit{ikke} 移到 \textit{kan} 前面）
  \item \textit{\da{Jeg er kun delvist enig i, at} onlinemøder altid er spild af tid.}（用温和的部分同意替代全盘否定，并把 \textit{med} 改为 \textit{i}）
\end{enumerate}

\textbf{练习 C 答案}
\begin{enumerate}
  \item \textit{Jeg er delvist enig med dig, men vi har stadig brug for flere tal.}
  \item \textit{Ganske vist er denne plan meget dyr, men ikke desto mindre har den stor langsigtet værdi.}
  \item \textit{Det kan godt være, at han har ret, men jeg synes stadig, vi bør vente.}
  \item \textit{Selvom alle er uenige, bevarer teamet dog en høflig og saglig tone.}
  \item \textit{Jeg er uenig med min kollega i dette spørgsmål om den digitale løsning.}
  \item \textit{Selvom onlinemøder sparer tid, kan de ikke helt erstatte den fysiske samtale.}
\end{enumerate}

\textbf{DU3 Modul 4 高分范文（Debatindlæg · 135 词）}\\
\textit{\textbf{Storrumskontorer kan fungere med klare rammer}}

\textit{Jeg forstår godt læserens frustration over støj på arbejdspladsen, og jeg er delvist enig i, at konstant uro kan ødelægge koncentrationen. \da{Det kan godt være, at} lukkede enmandskontorer giver mere ro ved skrivebordet, \dabold{men} et totalt forbud mod åbne kontorer er efter min mening ikke den rigtige vej.}

\textit{\da{Ganske vist} kan telefonopkald forstyrre kollegerne, \dabold{men ikke desto mindre} spiller åbne kontorer en vigtig rolle for samarbejde og hurtig sparring i hverdagen. Når man sidder sammen, bliver det langt lettere at stille et kort spørgsmål og hjælpe nye medarbejdere ind i opgaverne. \da{Selvom} støjen er en reel udfordring, \dabold{bør} vi \dabold{dog} vælge bedre indretning frem for at bygge vægge mellem alle.}

\textit{Med faste stillerum, tydelige samtaleregler og fleksible hjemmearbejdsdage kan vi nemlig bevare både roen og fællesskabet på moderne arbejdspladser.}

\small\textbf{名师点评}：正文含标题共 135 词，完全超越 Modul 4 Opgave 2 的 100 词基准。文章精准示范了丹麦式“礼貌反驳三步走”：首段先共情对手痛点并用 \textit{Det kan godt være, at..., men...} 温和转向；中段连用 \textit{Ganske vist..., men ikke desto mindre...}（双倒装满分无误）与 \textit{Selvom..., bør vi dog...}（不加 \textit{men}）展开反论；末段以建设性折中方案收束，逻辑严密、语调得体，达到典型 B2- 优异水准。
\end{answerbox}